\lecture{7}{2026-09-23}{The Principles of Analytic Combinatorics}

The asymptotic technique introduced in the previous lecture generalizes far beyond the tangent function.
If \(F(z) = \sum_{n \ge 0} f_n z^n\) is analytic at the origin, the Cauchy integral formula yields
\[
    f_n = \frac{1}{2\pi i} \int_{|z| = r} \frac{F(z)}{z^{n+1}} \, dz
\]
for any \(r > 0\) strictly less than the radius of convergence of \(F\).
Applying the Maximum Modulus Bound:
\[
    |f_n| \le \frac{1}{2\pi} \operatorname{length}(\{|z| = r\}) \cdot \max_{|z| = r} \left| \frac{F(z)}{z^{n+1}} \right|
    \le \frac{1}{2\pi} (2\pi r) \frac{M}{r^{n+1}} = \frac{M}{r^n},
\]
where \(M = \max_{|z| = r} |F(z)|\).
Consequently, \((f_n)_{n \ge 0}\) grows at most exponentially.

\section*{Exponential growth and dominant singularities}

\begin{definition}[Exponential growth]
    The \emph{exponential growth} of a sequence \((f_n)_{n \ge 0}\) is defined by
    \[
        \rho = \limsup_{n \to \infty} |f_n|^{1/n}.
    \]
\end{definition}

For example, if \(f_n = 2^n\), then \(\rho = 2\).
If \(f_n = 2^n(1 + (-1)^n)\), then \(\rho = 2\) as well.

By the Cauchy--Hadamard theorem (root test), \(\rho = \frac{1}{R}\), where \(R\) is the radius of convergence of \(\sum_{n \ge 0} f_n z^n\).
Because a power series converges up to its nearest singularity, this leads to the fundamental connection between analysis and asymptotics:

\begin{theorem}
    Let \(F(z) = \sum_{n \ge 0} f_n z^n\) have radius of convergence \(R > 0\).
    Then \(R\) equals the minimal modulus of a singularity of \(F(z)\), so
    \[
        \rho = \frac{1}{R}.
    \]
\end{theorem}

\begin{definition}[Dominant singularity]
    A singularity of \(F(z)\) of minimal modulus \(R\) is called a \emph{dominant singularity}.
\end{definition}

\begin{example}[Binary trees]
    Recall that the generating function of binary trees is \(B(z) = \frac{1 - \sqrt{1 - 4z}}{2z}\).
    The dominant singularity of \(B(z)\) is the branch point at \(z = 1/4\), so the exponential growth of binary trees is \(\rho = 4\).
    Therefore, \(b_n = 4^n \cdot \Lambda(n)\), where \(\Lambda(n)\) grows subexponentially.
\end{example}


\section*{The principles of analytic combinatorics}

These observations form the core philosophy of analytic combinatorics, formulated by Flajolet and Sedgewick:

\begin{itemize}
    \item \textbf{First Principle of Analytic Combinatorics}: The \emph{location} of a generating function's singularities determines the \emph{exponential growth} \(\rho\) of its coefficients.
    \item \textbf{Second Principle of Analytic Combinatorics}: The \emph{nature} (or \emph{type}) of a generating function's singularities determines the \emph{subexponential growth} of its coefficients.
\end{itemize}

\begin{example}
    Consider three generating functions, each with a singularity at \(z = 1/4\):
    \begin{itemize}
        \item If \(F(z) = \frac{1}{1 - 4z}\) (simple pole), then \(f_n = 4^n\).
        \item If \(F(z) = \frac{1}{(1 - 4z)^2}\) (double pole), then \(f_n = (n + 1) 4^n\).
        \item If \(F(z) = \frac{1 - \sqrt{1 - 4z}}{2z}\) (algebraic branch point), then \(f_n \sim \frac{4^n}{\sqrt{\pi n^3}}\).
    \end{itemize}
    In all three cases, the location \(z = 1/4\) fixes the base \(4^n\), while the nature of the singularity determines the subexponential factor (\(1\), \(n+1\), or \(n^{-3/2}/\sqrt{\pi}\)).
\end{example}


\section*{Asymptotics from polar singularities}

When the dominant singularities are poles, sequence asymptotics can be derived systematically by residue calculus.

\begin{proposition}
    \label{prop:rational-asymptotics}
    Suppose \(F(z) = \sum_{n \ge 0} f_n z^n\) is analytic on the circle \(\{|z| = R\}\) and analytic in the disk \(\{|z| < R\}\) except at a finite number of distinct nonzero poles \(\sigma_1, \dots, \sigma_m\).
    Then there exist polynomials \(P_1, \dots, P_m\) such that
    \[
        f_n = \sum_{k=1}^m P_k(n) \sigma_k^{-n} + O(R^{-n}),
    \]
    where the degree of \(P_k\) is one less than the order of the pole at \(\sigma_k\).
\end{proposition}

\begin{proof}
    Integrating over the circle \(|z| = R\) encloses the origin and the poles \(\sigma_1, \dots, \sigma_m\).
    By the Cauchy residue theorem,
    \[
        \frac{1}{2\pi i} \int_{|z| = R} \frac{F(z)}{z^{n+1}} \, dz = \Res_{z = 0} \left( \frac{F(z)}{z^{n+1}} \right) + \sum_{k=1}^m \Res_{z = \sigma_k} \left( \frac{F(z)}{z^{n+1}} \right).
    \]
    Since \(\Res_{z=0} \left( \frac{F(z)}{z^{n+1}} \right) = f_n\), rearranging gives
    \[
        f_n = - \sum_{k=1}^m \Res_{z = \sigma_k} \left( \frac{F(z)}{z^{n+1}} \right) + \frac{1}{2\pi i} \int_{|z| = R} \frac{F(z)}{z^{n+1}} \, dz.
    \]
    By the Maximum Modulus Bound, the integral is \(O(R^{-n})\).
    If \(\sigma_k\) is a pole of order \(d\), then
    \[
        \Res_{z = \sigma_k} \left( \frac{F(z)}{z^{n+1}} \right) = \frac{1}{(d-1)!} \lim_{z \to \sigma_k} \frac{d^{d-1}}{dz^{d-1}} \left( (z - \sigma_k)^d F(z) z^{-(n+1)} \right).
    \]
    By Leibniz's rule, differentiating \(z^{-(n+1)}\) up to \(d-1\) times produces factors of the form \((-1)^j (n+1)\dots(n+j) \sigma_k^{-(n+1+j)}\), which is a polynomial in \(n\) of degree \(d-1\) multiplied by \(\sigma_k^{-n}\).
    Absorbing the negative sign into \(P_k(n) = -\Res_{z=\sigma_k} \left( \frac{F(z)}{z^{n+1}} \right) \sigma_k^n\) yields the claim.
\end{proof}

\begin{example}[Derangements]
    The exponential generating function for derangements is \(D(z) = \sum_{n \ge 0} d_n \frac{z^n}{n!} = \frac{e^{-z}}{1 - z}\).
    Its only singularity in \(\mathbb{C}\) is a simple pole at \(z = 1\).
    Applying \Cref{prop:rational-asymptotics} for any choice of radius \(R > 1\):
    \[
        \frac{d_n}{n!} = - \Res_{z = 1} \left( \frac{e^{-z}}{z^{n+1}(1 - z)} \right) + O(R^{-n}).
    \]
    Since \(z = 1\) is a simple pole, \(\Res_{z=1} \frac{e^{-z}}{z^{n+1}(1-z)} = \frac{e^{-1}}{1^{n+1}(-1)} = -e^{-1}\).
    Therefore,
    \[
        \frac{d_n}{n!} = e^{-1} + O(R^{-n}) \to \frac{1}{e} \text{as } n \to \infty.
    \]
    Thus, the probability that a uniformly random permutation of size \(n\) has no fixed points rapidly approaches \(1/e\), with exponential convergence for any rate \(R > 1\).
\end{example}